% ==========================================================================
% Capítulo: Estado del arte
% Estado: se registra en ../STATUS.md, no aquí.
%
% Debe contener (project-context/formato-anteproyecto.md, "Estado del
% arte"; extensión máxima 4 páginas):
%   - Revisión documental del conocimiento acumulado alrededor del
%     problema: ¿cómo se ha resuelto antes?
%   - Alternativas de solución previas, caracterizadas y analizadas, con
%     ventajas y desventajas de cada una.
%   - Argumentación explícita de por qué esas propuestas NO resuelven el
%     problema en el contexto de IAsLab / Universidad Icesi.
% NO debe contener: teoría de base (eso es 05-marco-teorico.tex) ni
% literatura gris sin filtro. Fuentes: IEEE, ACM, libros acreditados.
%
% Candidata ya identificada por los autores (2026-09-12), pendiente de
% investigar con fuente citable real, NO redactar todavía a partir de esta
% nota: las plataformas MLOps gestionadas en la nube (TrueFoundry, Amazon
% SageMaker, Microsoft Azure Machine Learning, Google Vertex AI) resuelven
% entrenamiento gestionado, registro de modelos, pipelines y despliegue con
% autoescalado, pero no sirven para el problema del IAsLab por una razón
% estructural, no técnica: son servicios en la nube, y el proyecto exige
% operar estrictamente on-premise sobre el hardware ya instalado en la sala
% 104M (restricción confirmada en project-context/documentation.md, sección
% Scope). El Investigador debe traer una fuente real (documentación oficial
% de cada plataforma, o una comparativa académica) antes de que esto entre
% como afirmación citada — ver skills/reference-writting/recency.md para el
% corte de antigüedad (2015 en adelante) también aplicado aquí.
% ==========================================================================
\chapter{Estado del arte}\label{cha:estado-del-arte}

% [CONTEXTO Y GUÍA COMPLETA PARA REDACTAR EL ESTADO DEL ARTE]
% 
% Bienvenido(a) a la guía definitiva para elaborar el estado del arte, si eres estudiante o investigador, sabes que esta sección es crucial en cualquier trabajo académico o científico, ya que te permite conocer y analizar las fuentes existentes sobre un tema de investigación.
% 
% El estado del arte sirve al investigador como referencia para adoptar una postura crítica sobre lo que se ha hecho y lo que aún falta por hacer en relación con una temática o problemática específica. La búsqueda de evidencias, es decir, referencias a trabajos similares, se convierte en el punto de partida del estado del arte.
% 
% Guía APA aquí: https://normasapa.in/estado-del-arte/
% 
% ¿Qué es el estado del arte?
% 
% El estado del arte es un tipo de investigación documental que consiste en examinar la literatura disponible sobre un tema específico. En palabras más sencillas, el estado del arte es un proceso de investigación que implica leer y analizar diferentes textos académicos para conocer lo que se ha investigado sobre un tema en particular. Al hacer un estado del arte, aprendes sobre los avances, desafíos y tendencias de la investigación en ese tema. De esta manera, puedes desarrollar un conocimiento crítico basado en la revisión y análisis de distintos tipos de textos.Lee Revistas Científicas
% 
% Es importante destacar que un estado del arte no es simplemente una lista de textos con resúmenes cortos de cada uno, sino que implica la selección y análisis de las partes más relevantes para la investigación. Además, se debe establecer un sistema de parámetros de descarte e inclusión de los textos consultados para garantizar la sistemática y replicabilidad del proceso.
% Estado del arte: Mapa mental
% Estado del arte: Mapa mental
% 
% Guía APA aquí: https://normasapa.in/estado-del-arte/
% 
% ¿Cómo escribir el estado del arte?
% 
% El estado del arte permite dar validez a cualquier investigación; esta parte del trabajo proporciona una visión general del conocimiento existente sobre un tema específico le sirve al indagador como referencia. Sigue los siguientes pasos (Junto con ejemplos) para incluir el estado del arte en cualquier trabajo escrito:
% Identificar el tema de investigación
% 
% El primer paso para hacer el estado del arte es identificar el tema de investigación. El tema debe ser específico y estar claramente definido. Esto permitirá al investigador enfocarse en la literatura relevante y evitar la revisión de información no pertinente. El investigador también debe tener un conocimiento previo sobre el tema, ya que esto puede facilitar la identificación de lo relevante.
% Por ejemplo: Supongamos que estás interesado en investigar sobre el efecto de la música en la memoria a corto plazo. El primer paso sería identificar el tema de investigación a fondo; en pocas palabras, se debe conocer o buscar información relacionados con la música y la memoria a corto plazo.Citas bibliográficas guía
% Definir la estrategia de búsqueda
% 
% La estrategia de búsqueda debe incluir palabras clave relevantes y criterios de inclusión y exclusión. Las palabras clave deben ser específicas y relevantes para el tema de investigación y deben cubrir todas las posibles variaciones de la terminología utilizada en el campo.
% 
%     Los criterios de inclusión y exclusión deben ser claros y deben definir claramente los tipos de documentos que se incluirán y excluirán de la revisión bibliográfica.
% 
% Por ejemplo: Para definir la estrategia de búsqueda implica determinar las palabras clave y los términos de búsqueda que se utilizarán para encontrar la información relevante. Por ejemplo, en nuestro ejemplo, se podrían utilizar términos como «música», «memoria a corto plazo», «estudios de investigación», etc.
% 
% Guía APA aquí: https://normasapa.in/estado-del-arte/
% 
% Identificar las fuentes de información
% 
% El investigador debe ser lo más exhaustivo posible al identificar las fuentes relevantes. Recuerda que no todas las fuentes de información son iguales y que algunas pueden ser más confiables que otras.
% 
%     Puede que hayan 1000, incluso más investigaciones que traten sobre un tema; para evitar aburrimiento al lector y hojas sin leer, incluye únicamente aquellas relevantes o que aporten y/o atraigan.
% 
% Por ejemplo: Es importante identificar las fuentes de información relevantes. Estas pueden incluir bases de datos académicas, revistas científicas, libros, tesis doctorales, entre otros. Algunas de las fuentes de información relevantes para este tema podrían ser la base de datos de PsycINFO, la revista Memory and Cognition y el libro «The Cognitive Neuroscience of Music».
% Realizar la revisión bibliográfica
% 
% Se deben revisar todos los documentos identificados para determinar su relevancia para el tema de investigación. La revisión bibliográfica debe incluir la lectura completa de los documentos relevantes y la extracción de información relevante. Intenta resumir la información extraída pues cada fuente debe incluir los objetivos, la metodología utilizada, los resultados y las conclusiones.Marco teórico guía
% 
%     Si en una búsqueda encontraste 1000 fuentes y del filtrado sacaste 5 o 10 confiables, verificables y que respondan el punto de inicio; debes considerarlo como un excelente inicio.
% 
% Por ejemplo: En este caso, se podría leer el estudio «Effects of Music on Memory for Text» de Wallace y otros (1994) para obtener información sobre el efecto de la música en la memoria a corto plazo. Para éste punto es necesario que el investigador lea y analice los materiales relevantes para el tema.
% 
% Guía APA aquí: https://normasapa.in/estado-del-arte/
% 
% Sintetizar la información
% 
% Una vez que se han revisado todos los documentos relevantes, el siguiente paso es sintetizar la información. La información extraída de los documentos relevantes debe agruparse según temas y subtemas. El investigador debe determinar las tendencias actuales y los vacíos en la investigación.
% 
%     La información sintetizada debe ser lo más clara posible y debe permite comprender fácilmente el estado actual del conocimiento en el campo.
%     Esto puede implicar resumir los resultados de los estudios de investigación relevantes, identificar tendencias en el campo y discutir los hallazgos clave.
% 
% Por ejemplo: Se podría sintetizar la información encontrada en los estudios de investigación relevantes para identificar los efectos de la música en la memoria a corto plazo y discutir las tendencias en el campo.
% Escribir el estado del arte
% 
% El siguiente paso es escribir el estado del arte; está sección debe incluir una introducción al tema de investigación y una breve descripción de la metodología utilizada para realizar la revisión bibliográfica.Consulta Enciclopedias Digitales
% 
%     La sección principal debe presentar los resultados de la revisión bibliográfica. Los resultados deben presentarse de manera lógica y organizada, utilizando subtemas y secciones.
%     Es importante que el estado del arte sea lo más claro y conciso posible, utilizando un lenguaje técnico y científico adecuado para el tema de investigación.
% 
% Por ejemplo: En este caso, se podría escribir un estado del arte que incluya una introducción al tema de la música y la memoria a corto plazo, una descripción de la estrategia de búsqueda utilizada y los resultados de la revisión bibliográfica.
% 
% Guía APA aquí: https://normasapa.in/estado-del-arte/
% 
% Revisar y actualizar el estado del arte
% 
% Una vez que se ha escrito el estado del arte, el siguiente paso es revisarlo y actualizarlo regularmente. El estado del arte debe revisarse antes de la presentación del trabajo de investigación final para asegurarse de que toda la información sea relevante y esté actualizada.
% Por ejemplo: Si se encuentra un nuevo estudio de investigación relevante para el tema, se deberá agregar la información pertinente al estado del arte existente y actualizar la lista de referencias.
% ¿Cuál es la estructura del estado del arte?
% 
% El estado del arte no es obligatorio de incluir en los trabajos escritos, es decisión de la institución o educador el que lo debas poner o no; si éste es tu caso debes tener en cuenta que la estructura puede variar según el campo de estudio, pero por norma general debe incluir la siguiente información:Asesoría bibliográfica
% Introducción
% 
% La introducción debe presentar una breve presentación del tema a tratar y su importancia en el campo de estudio. Además de la relevancia, antecedentes, objetivos generales del problema a tratar; es recomendable responder las siguientes preguntas para contextualizar al lector:
% 
%     ¿Por qué he decidido indagar sobre éste problema?
%     ¿Cómo se definieron los problemas?
%     ¿Qué contenidos, tópicos o dimensiones, se han definido como prioritarios?
%     ¿Cuáles son las características principales del tema?
%     ¿Es relevante el problema/tema elegido?
% 
% Guía APA aquí: https://normasapa.in/estado-del-arte/
% 
% Investigaciones (Antecedentes)
% 
% Un trabajo escrito sin otras fuentes de investigación NO es confiable ante los ojos del lector o cualquier otro investigador. Los antecedentes deben incluir una revisión de la literatura existente sobre el tema, incluyendo trabajos previos relevantes y teorías relacionadas.
% 
% El estado del arte es una parte del proyecto que presenta el punto de vista y resultados obtenidos por otros indagadores, eso si, con relación con el tema a tratar. Incluye fuentes, resultados y problemas:
% 
% Guía APA aquí: https://normasapa.in/estado-del-arte/
% 
% Debes agregar referencias directas a estás investigaciones agregando una descripción de los resultados obtenidos; recuerda dar orden a cada fuente agregada, puede ser: Cronológico, geográfico o por palabras claves; para contextualizar más puedes agregar los siguientes elementos:
% 
%     ¿Qué evidencias empíricas y metodológicas se utilizaron?
%     Una descripción de los métodos utilizados para llevar a cabo la revisión de los datos.
%     Una síntesis de los resultados de la revisión de la literatura, identificando los avances más significativos en el campo y las áreas donde se necesita más investigación.
%     ¿Qué temas de indagación se han definido como relacionados con el tema de la investigación?
%     Marco teórico guía
% 
% Conclusiones
% 
% Describe los resultados, descubrimientos, discusiones, limitaciones, problemas indagados, dudas sin aclarar. Un buen párrafo de conclusión debe responder las siguientes preguntas e incluir:
% 
%     ¿Cuál es el resultado esperado de las investigaciones?
%     ¿Cuántas limitaciones hubieron para la búsqueda de resultados?
%     ¿Cuáles fueron las fortalezas y debilidades de los estudios revisados?
%     ¿Qué problemas quedan por resolver o salieron en el proceso?
%     ¿Fueron los resultados esperados?
% 
% Guía APA aquí: https://normasapa.in/estado-del-arte/
% 
% Inicio
% ›
% Estructura
% › Guía completa para escribir un Estado del Arte, con ejemplos prácticos
% Guía completa para escribir un Estado del Arte, con ejemplos prácticos
% marzo 20, 2023 por Andrés Rivas	
% 
% Bienvenido(a) a la guía definitiva para elaborar el estado del arte, si eres estudiante o investigador, sabes que esta sección es crucial en cualquier trabajo académico o científico, ya que te permite conocer y analizar las fuentes existentes sobre un tema de investigación.
% 
% El estado del arte sirve al investigador como referencia para adoptar una postura crítica sobre lo que se ha hecho y lo que aún falta por hacer en relación con una temática o problemática específica. La búsqueda de evidencias, es decir, referencias a trabajos similares, se convierte en el punto de partida del estado del arte.
% 
% También te recomendamos: Cómo escribir el marco teórico en una investigación
% 
% ÍNDICE DE CONTENIDOS (Omitido por brevedad en los comentarios, referirse a la guía completa arriba)
% 
% ¿Cómo realizar la búsqueda de antecedentes para escribir el estado del arte?
% 
% Es importante tener en cuenta que la búsqueda de fuentes para un estado del arte es un proceso iterativo y continuo, por lo que es recomendable revisar y actualizar la búsqueda de manera regular.Lee Revistas Científicas
% 
%     Define el tema: Es importante tener claridad sobre el tema o campo de estudio que se va a investigar, ya que esto ayudará a enfocar la búsqueda de fuentes de información.
%     Identifica las palabras clave: Estos términos deben ser relevantes al tema de investigación.
%     Utiliza bases de datos académicas: Existen diversas bases de datos académicas que contienen una gran cantidad de información científica y técnica, tales como Scopus, Web of Science, PubMed, etc. Estas bases de datos permiten realizar búsquedas avanzadas utilizando diferentes criterios, tales como palabras clave, fecha de publicación, tipo de documento, entre otros.
%     Revisa revistas científicas: Aquellas especializadas en el área de estudio son una excelente fuente de información. Es necesario revisar los últimos números de las revistas y buscar artículos relevantes al tema de investigación.
%     Consulta bibliografías de artículos y libros: Las bibliografías en cualquier tipo de trabajo escrito pueden ser una excelente fuente para encontrar referencias adicionales y ampliar la búsqueda.
%     Utiliza herramientas de búsqueda en línea: Las herramientas en línea como Google Académico, ResearchGate, Academia.edu, entre otras, pueden ser útiles para encontrar información adicional.
%     Consulta a expertos en el tema: Es importante consultar a expertos en el área de estudio para obtener información adicional y recomendaciones de fuentes de información relevantes.
% 
% Guía APA aquí: https://normasapa.in/estado-del-arte/
% 
% ¿Porqué es necesario incluir el estado del arte en una investigación?
% 
% El estado del arte es importante porque proporciona contexto, identifica vacíos, valida la relevancia y originalidad del trabajo, identifica metodologías y técnicas, y proporciona una lista de referencias útiles para otros investigadores o futuras investigaciones del tema a tratar; además:
% 
%     Contextualiza: El estado del arte proporciona contexto sobre el tema que se está tratando en el trabajo. Esto permite a los lectores comprender la evolución histórica del tema y las tendencias actuales.
%     Identifica los vacíos: Estos vacíos pueden ser oportunidades para futuras investigaciones y pueden ser el punto de partida para nuevos proyectos.
%     Valida la relevancia: Esto ayuda a los autores a demostrar que su trabajo es una contribución significativa al conocimiento existente.
%     Brinda metodologías: El estado del arte también puede proporcionar información sobre las metodologías y técnicas que se han utilizado en investigaciones anteriores sobre el tema. Esto puede ayudar a los autores a seleccionar las metodologías y técnicas adecuadas para su trabajo.
%      Referencias bibliográficas: El estado del arte también puede servir como una lista de referencias bibliográficas útiles para otros investigadores interesados en el tema. Esto ayuda a establecer una red de conocimientos y fomenta la colaboración entre investigadores.
% 
% ¿Cuáles son los objetivos de un estado del arte?
% 
% Como mencionamos al principio de este artículo, el estado del arte es un tipo de investigación documental acerca de la forma en que diferentes autores han tratado un tema específico; si aún no te queda clara su importancia, te recomendamos leer los siguientes objetivos:Consulta Enciclopedias Digitales
% 
% Guía APA aquí: https://normasapa.in/estado-del-arte/
% 
% Amplia el conocimiento: El estado del arte es la investigación de otra investigación, es decir, cada proyecto saliente de un mismo tema ofrecerá distintas hipótesis, problemas, resultados o conclusiones
% Fuente de referencias: Para que una investigación sea confiable, es necesario aportar referencias del mismo tema/área. Un proyecto, tesis o trabajo escrito con estado del arte debe adjuntar todas las fuentes necesarias y esto será un aporte de investigación para el próximo indagador
% Relevancia: Cuando existe mucha información, referencias o indagadores en un mismo tema, eso demuestra el grado de importancia/relevancia del tema
% Diferentes puntos de vista: No todos los escritores llegarán a la misma conclusión del tema o problema indagado, muchos ni siquiera tendrán resultados y lanzarán críticas o diferentes puntos de vista; al final, el propósito de toda investigación es ofrecer algo más al lector
% Plantea otros problemas: Muchas investigaciones hoy en día tienen problemas de un mismo tema que aún están sin resolver. Incluir el estado del arte le permite al lector conocer las limitaciones y dudas sin aclarar; además de aquellas preguntas que surgen al realizar la investigación
% Al realizar el estado del arte, el indagador/estudiante consolida un conocimiento crítico basado en la lectura y el análisis de diferentes tipos de textos
% [verify: not yet written]
